\begin{abstract}
Active imaging with diffusion priors is commonly driven by disagreement
between endpoint samples or by a single local covariance.  The former spends
many denoiser evaluations, while the latter Gaussianizes away the outcome
dependence that makes adaptive planning useful.  We introduce \method, a
compressed belief state built from a small set of measurement-guided diffusion
branches, each augmented by a low-rank local Tweedie/Jacobian factor.  Under
this finite-mixture surrogate, mutual information decomposes into an analytic
within-branch term and a bounded mode-disambiguation term.  More importantly,
all action likelihoods, information bounds, Bayesian updates, and fantasy
outcomes live in at most $\sum_k r_k+K-1$ whitened coordinates, independent of
the scalar size of an MRI line or LiDAR angular tile.  This enables
outcome-contingent horizon-two planning without additional denoiser calls and
an analytic correction that is fed back into reverse diffusion.  We also
derive a target-information extension that marginalizes low-dimensional
forward-operator nuisance instead of treating calibration as known.  Synthetic
falsification tests reduce covariance-estimation MSE by $61.6\%$, compress a
4,096-dimensional action to 11 coordinates, and recover the correct
non-myopic router.  The locked imaging study evaluates raw multicoil
fastMRI+ brain/knee data and KITTI-360 programmable LiDAR at matched
measurement and denoiser budgets.  Imaging results in this working draft are
explicitly \pending{} and will not be claimed before the preregistered
patient/drive-level gates pass.
\end{abstract}
